\IfFileExists{stacks-project.cls}{%
\documentclass{stacks-project}
}{%
\documentclass{amsart}
}

% For dealing with references we use the comment environment
\usepackage{verbatim}
\newenvironment{reference}{\comment}{\endcomment}
%\newenvironment{reference}{}{}
\newenvironment{slogan}{\comment}{\endcomment}
\newenvironment{history}{\comment}{\endcomment}

% For commutative diagrams we use Xy-pic
\usepackage[all]{xy}

% We use 2cell for 2-commutative diagrams.
\xyoption{2cell}
\UseAllTwocells

% We use multicol for the list of chapters between chapters
\usepackage{multicol}

% This is generally recommended for better output
\usepackage{lmodern}
\usepackage[T1]{fontenc}

% For cross-file-references
\usepackage{xr-hyper}

% Package for hypertext links:
\usepackage{hyperref}

% For any local file, say "hello.tex" you want to link to please
% use \externaldocument[hello-]{hello}
\externaldocument[introduction-]{introduction}
\externaldocument[conventions-]{conventions}
\externaldocument[sets-]{sets}
\externaldocument[categories-]{categories}
\externaldocument[topology-]{topology}
\externaldocument[sheaves-]{sheaves}
\externaldocument[sites-]{sites}
\externaldocument[stacks-]{stacks}
\externaldocument[fields-]{fields}
\externaldocument[algebra-]{algebra}
\externaldocument[brauer-]{brauer}
\externaldocument[homology-]{homology}
\externaldocument[derived-]{derived}
\externaldocument[simplicial-]{simplicial}
\externaldocument[more-algebra-]{more-algebra}
\externaldocument[smoothing-]{smoothing}
\externaldocument[modules-]{modules}
\externaldocument[sites-modules-]{sites-modules}
\externaldocument[injectives-]{injectives}
\externaldocument[cohomology-]{cohomology}
\externaldocument[sites-cohomology-]{sites-cohomology}
\externaldocument[dga-]{dga}
\externaldocument[dpa-]{dpa}
\externaldocument[sdga-]{sdga}
\externaldocument[hypercovering-]{hypercovering}
\externaldocument[schemes-]{schemes}
\externaldocument[constructions-]{constructions}
\externaldocument[properties-]{properties}
\externaldocument[morphisms-]{morphisms}
\externaldocument[coherent-]{coherent}
\externaldocument[divisors-]{divisors}
\externaldocument[limits-]{limits}
\externaldocument[varieties-]{varieties}
\externaldocument[topologies-]{topologies}
\externaldocument[descent-]{descent}
\externaldocument[perfect-]{perfect}
\externaldocument[more-morphisms-]{more-morphisms}
\externaldocument[flat-]{flat}
\externaldocument[groupoids-]{groupoids}
\externaldocument[more-groupoids-]{more-groupoids}
\externaldocument[etale-]{etale}
\externaldocument[chow-]{chow}
\externaldocument[intersection-]{intersection}
\externaldocument[pic-]{pic}
\externaldocument[weil-]{weil}
\externaldocument[adequate-]{adequate}
\externaldocument[dualizing-]{dualizing}
\externaldocument[duality-]{duality}
\externaldocument[discriminant-]{discriminant}
\externaldocument[derham-]{derham}
\externaldocument[local-cohomology-]{local-cohomology}
\externaldocument[algebraization-]{algebraization}
\externaldocument[curves-]{curves}
\externaldocument[resolve-]{resolve}
\externaldocument[models-]{models}
\externaldocument[functors-]{functors}
\externaldocument[equiv-]{equiv}
\externaldocument[pione-]{pione}
\externaldocument[etale-cohomology-]{etale-cohomology}
\externaldocument[proetale-]{proetale}
\externaldocument[relative-cycles-]{relative-cycles}
\externaldocument[more-etale-]{more-etale}
\externaldocument[trace-]{trace}
\externaldocument[crystalline-]{crystalline}
\externaldocument[spaces-]{spaces}
\externaldocument[spaces-properties-]{spaces-properties}
\externaldocument[spaces-morphisms-]{spaces-morphisms}
\externaldocument[decent-spaces-]{decent-spaces}
\externaldocument[spaces-cohomology-]{spaces-cohomology}
\externaldocument[spaces-limits-]{spaces-limits}
\externaldocument[spaces-divisors-]{spaces-divisors}
\externaldocument[spaces-over-fields-]{spaces-over-fields}
\externaldocument[spaces-topologies-]{spaces-topologies}
\externaldocument[spaces-descent-]{spaces-descent}
\externaldocument[spaces-perfect-]{spaces-perfect}
\externaldocument[spaces-more-morphisms-]{spaces-more-morphisms}
\externaldocument[spaces-flat-]{spaces-flat}
\externaldocument[spaces-groupoids-]{spaces-groupoids}
\externaldocument[spaces-more-groupoids-]{spaces-more-groupoids}
\externaldocument[bootstrap-]{bootstrap}
\externaldocument[spaces-pushouts-]{spaces-pushouts}
\externaldocument[spaces-chow-]{spaces-chow}
\externaldocument[groupoids-quotients-]{groupoids-quotients}
\externaldocument[spaces-more-cohomology-]{spaces-more-cohomology}
\externaldocument[spaces-simplicial-]{spaces-simplicial}
\externaldocument[spaces-duality-]{spaces-duality}
\externaldocument[formal-spaces-]{formal-spaces}
\externaldocument[restricted-]{restricted}
\externaldocument[spaces-resolve-]{spaces-resolve}
\externaldocument[formal-defos-]{formal-defos}
\externaldocument[defos-]{defos}
\externaldocument[cotangent-]{cotangent}
\externaldocument[examples-defos-]{examples-defos}
\externaldocument[algebraic-]{algebraic}
\externaldocument[examples-stacks-]{examples-stacks}
\externaldocument[stacks-sheaves-]{stacks-sheaves}
\externaldocument[criteria-]{criteria}
\externaldocument[artin-]{artin}
\externaldocument[quot-]{quot}
\externaldocument[stacks-properties-]{stacks-properties}
\externaldocument[stacks-morphisms-]{stacks-morphisms}
\externaldocument[stacks-limits-]{stacks-limits}
\externaldocument[stacks-cohomology-]{stacks-cohomology}
\externaldocument[stacks-perfect-]{stacks-perfect}
\externaldocument[stacks-introduction-]{stacks-introduction}
\externaldocument[stacks-more-morphisms-]{stacks-more-morphisms}
\externaldocument[stacks-geometry-]{stacks-geometry}
\externaldocument[moduli-]{moduli}
\externaldocument[moduli-curves-]{moduli-curves}
\externaldocument[examples-]{examples}
\externaldocument[exercises-]{exercises}
\externaldocument[guide-]{guide}
\externaldocument[desirables-]{desirables}
\externaldocument[coding-]{coding}
\externaldocument[obsolete-]{obsolete}
\externaldocument[fdl-]{fdl}
\externaldocument[index-]{index}

% Theorem environments.
%
\theoremstyle{plain}
\newtheorem{theorem}[subsection]{Theorem}
\newtheorem{proposition}[subsection]{Proposition}
\newtheorem{lemma}[subsection]{Lemma}

\theoremstyle{definition}
\newtheorem{definition}[subsection]{Definition}
\newtheorem{example}[subsection]{Example}
\newtheorem{exercise}[subsection]{Exercise}
\newtheorem{situation}[subsection]{Situation}

\theoremstyle{remark}
\newtheorem{remark}[subsection]{Remark}
\newtheorem{remarks}[subsection]{Remarks}

\numberwithin{equation}{subsection}

% Macros
%
\def\lim{\mathop{\mathrm{lim}}\nolimits}
\def\colim{\mathop{\mathrm{colim}}\nolimits}
\def\Spec{\mathop{\mathrm{Spec}}}
\def\Hom{\mathop{\mathrm{Hom}}\nolimits}
\def\Ext{\mathop{\mathrm{Ext}}\nolimits}
\def\SheafHom{\mathop{\mathcal{H}\!\mathit{om}}\nolimits}
\def\SheafExt{\mathop{\mathcal{E}\!\mathit{xt}}\nolimits}
\def\Sch{\mathit{Sch}}
\def\Mor{\mathop{\mathrm{Mor}}\nolimits}
\def\Ob{\mathop{\mathrm{Ob}}\nolimits}
\def\Sh{\mathop{\mathit{Sh}}\nolimits}
\def\NL{\mathop{N\!L}\nolimits}
\def\CH{\mathop{\mathrm{CH}}\nolimits}
\def\proetale{{pro\text{-}\acute{e}tale}}
\def\etale{{\acute{e}tale}}
\def\QCoh{\mathit{QCoh}}
\def\Ker{\mathop{\mathrm{Ker}}}
\def\Im{\mathop{\mathrm{Im}}}
\def\Coker{\mathop{\mathrm{Coker}}}
\def\Coim{\mathop{\mathrm{Coim}}}

% Boxtimes
%
\DeclareMathSymbol{\boxtimes}{\mathbin}{AMSa}{"02}

%
% Macros for moduli stacks/spaces
%
\def\QCohstack{\mathcal{QC}\!\mathit{oh}}
\def\Cohstack{\mathcal{C}\!\mathit{oh}}
\def\Spacesstack{\mathcal{S}\!\mathit{paces}}
\def\Quotfunctor{\mathrm{Quot}}
\def\Hilbfunctor{\mathrm{Hilb}}
\def\Curvesstack{\mathcal{C}\!\mathit{urves}}
\def\Polarizedstack{\mathcal{P}\!\mathit{olarized}}
\def\Complexesstack{\mathcal{C}\!\mathit{omplexes}}
% \Pic is the operator that assigns to X its picard group, usage \Pic(X)
% \Picardstack_{X/B} denotes the Picard stack of X over B
% \Picardfunctor_{X/B} denotes the Picard functor of X over B
\def\Pic{\mathop{\mathrm{Pic}}\nolimits}
\def\Picardstack{\mathcal{P}\!\mathit{ic}}
\def\Picardfunctor{\mathrm{Pic}}
\def\Deformationcategory{\mathcal{D}\!\mathit{ef}}
\usepackage{cleveref}
\crefname{originalreference}{original reference}{original references}
\Crefname{originalreference}{Original reference}{Original references}
\makeatletter
\newlabel{categories-lemma-yoneda}{{001P}{source}{}{sourcepointer-1}{}}
\newlabel{categories-lemma-yoneda@cref}{{[originalreference][0][]001P}{[1][0][]source}{}{}{}}
\makeatother
\crefname{originaldefinition}{original definition}{original definitions}
\Crefname{originaldefinition}{Original definition}{Original definitions}
\makeatletter
\newlabel{definition-Serre-functor}{{0FY5}{source}{}{sourcepointer-2}{}}
\newlabel{definition-Serre-functor@cref}{{[originaldefinition][0][]0FY5}{[1][0][]source}{}{}{}}
\makeatother
\crefname{originalreference}{original reference}{original references}
\Crefname{originalreference}{Original reference}{Original references}
\makeatletter
\newlabel{derived-lemma-representable-homological}{{0149}{source}{}{sourcepointer-3}{}}
\newlabel{derived-lemma-representable-homological@cref}{{[originalreference][0][]0149}{[1][0][]source}{}{}{}}
\makeatother
\crefname{originalreference}{original reference}{original references}
\Crefname{originalreference}{Original reference}{Original references}
\makeatletter
\newlabel{homology-lemma-five-lemma}{{05QB}{source}{}{sourcepointer-4}{}}
\newlabel{homology-lemma-five-lemma@cref}{{[originalreference][0][]05QB}{[1][0][]source}{}{}{}}
\makeatother
\crefname{originalremark}{original remark}{original remarks}
\Crefname{originalremark}{Original remark}{Original remarks}
\makeatletter
\newlabel{remark-dual-matress}{{0H8C}{source}{}{sourcepointer-5}{}}
\newlabel{remark-dual-matress@cref}{{[originalremark][0][]0H8C}{[1][0][]source}{}{}{}}
\makeatother
\crefname{originalremark}{original remark}{original remarks}
\Crefname{originalremark}{Original remark}{Original remarks}
\makeatletter
\newlabel{remark-matress}{{0H8B}{source}{}{sourcepointer-6}{}}
\newlabel{remark-matress@cref}{{[originalremark][0][]0H8B}{[1][0][]source}{}{}{}}
\makeatother
\title{Uniqueness and exactness of Serre functors}
\author{The Stacks Project Authors}
\date{stacks-project:equiv.tex:local-snapshot-20261012; GNU FDL 1.2 or later}
\begin{document}
\maketitle
\noindent\textit{Allowed standard prerequisites (editorial).}\begin{itemize}
\item Yoneda lemma. categories-lemma-yoneda (cross-chapter, named). Identifies representing objects up to unique isomorphism.
\item Representable homological functors. derived-lemma-representable-homological (cross-chapter, named). Exactness of the dual Hom sequences.
\item Matress formalism. remark-matress, remark-dual-matress (same chapter, not carded). Constructs the comparison arrow phi between the two triangles.
\end{itemize}\medskip
\ifdefined\hypertarget\hypertarget{sourcepointer-1}{}\fi
\noindent\textit{Source reference.} Original reference 001P: focused-seeds/geometry/S9-eqv-Serre-functor/ref-tags.txt, source lines 5--5. Navigation pointer only: the counted proof cites this statement by \textbackslash{}ref as a standing input; the deduction the counted proof has to make is carried out inside the counted span, and no part of that deduction is delegated to this pointer.\par\medskip
\ifdefined\hypertarget\hypertarget{sourcepointer-2}{}\fi
\noindent\textit{Source reference.} Original definition 0FY5: sources/stacks-equiv/pinned-original-source.tex, source lines 378--378. Navigation pointer only: the counted proof cites this statement by \textbackslash{}ref as a standing input; the deduction the counted proof has to make is carried out inside the counted span, and no part of that deduction is delegated to this pointer.\par\medskip
\ifdefined\hypertarget\hypertarget{sourcepointer-3}{}\fi
\noindent\textit{Source reference.} Original reference 0149: focused-seeds/geometry/S9-eqv-Serre-functor/ref-tags.txt, source lines 6--6. Navigation pointer only: the counted proof cites this statement by \textbackslash{}ref as a standing input; the deduction the counted proof has to make is carried out inside the counted span, and no part of that deduction is delegated to this pointer.\par\medskip
\ifdefined\hypertarget\hypertarget{sourcepointer-4}{}\fi
\noindent\textit{Source reference.} Original reference 05QB: focused-seeds/geometry/S9-eqv-Serre-functor/ref-tags.txt, source lines 7--7. Navigation pointer only: the counted proof cites this statement by \textbackslash{}ref as a standing input; the deduction the counted proof has to make is carried out inside the counted span, and no part of that deduction is delegated to this pointer.\par\medskip
\ifdefined\hypertarget\hypertarget{sourcepointer-5}{}\fi
\noindent\textit{Source reference.} Original remark 0H8C: sources/stacks-equiv/pinned-original-source.tex, source lines 421--421. Navigation pointer only: the counted proof cites this statement by \textbackslash{}ref as a standing input; the deduction the counted proof has to make is carried out inside the counted span, and no part of that deduction is delegated to this pointer.\par\medskip
\ifdefined\hypertarget\hypertarget{sourcepointer-6}{}\fi
\noindent\textit{Source reference.} Original remark 0H8B: sources/stacks-equiv/pinned-original-source.tex, source lines 390--390. Navigation pointer only: the counted proof cites this statement by \textbackslash{}ref as a standing input; the deduction the counted proof has to make is carried out inside the counted span, and no part of that deduction is delegated to this pointer.\par\medskip
\begin{lemma}
\label{lemma-Serre-functor}
In the situation of Definition \ref{definition-Serre-functor}.
If a Serre functor exists, then it is unique up to unique isomorphism and
it is an exact functor of triangulated categories.
\end{lemma}

\begin{proof}
Given a Serre functor $S$ the object $S(X)$ represents
the functor $Y \mapsto \Hom_\mathcal{T}(X, Y)^\vee$.
Thus the object $S(X)$ together with the functorial identification
$\Hom_\mathcal{T}(X, Y)^\vee = \Hom_\mathcal{T}(Y, S(X))$
is determined up to unique isomorphism by the Yoneda lemma
(Categories, Lemma \ref{categories-lemma-yoneda}).
Moreover, for $\varphi : X \to X'$, the arrow $S(\varphi) : S(X) \to S(X')$
is uniquely determined by the corresponding transformation of functors
$\Hom_\mathcal{T}(X, -)^\vee \to \Hom_\mathcal{T}(X', -)^\vee$.

\medskip\noindent
For objects $X, Y$ of $\mathcal{T}$ we have
\begin{align*}
\Hom(Y, S(X)[1])^\vee
& =
\Hom(Y[-1], S(X))^\vee \\
& =
\Hom(X, Y[-1]) \\
& =
\Hom(X[1], Y) \\
& =
\Hom(Y, S(X[1]))^\vee
\end{align*}
By the Yoneda lemma we conclude that there is a unique isomorphism
$S(X[1]) \to S(X)[1]$ inducing the isomorphism from top left to bottom right.
Since each of the isomorphisms above is functorial in both $X$ and $Y$
we find that this defines an isomorphism of functors
$S \circ [1] \to [1] \circ S$.

\medskip\noindent
Let $(A, B, C, f, g, h)$ be a distinguished triangle in $\mathcal{T}$.
We have to show that the triangle $(S(A), S(B), S(C), S(f), S(g), S(h))$
is distinguished. Here we use the canonical isomorphism $S(A[1]) \to S(A)[1]$
constructed above to identify the target $S(A[1])$ of $S(h)$ with $S(A)[1]$.
We first observe that for any $X$ in $\mathcal{T}$
the triangle $(S(A), S(B), S(C), S(f), S(g), S(h))$ induces
a long exact sequence
$$
\ldots \to
\Hom(X, S(A)) \to
\Hom(X, S(B)) \to
\Hom(X, S(C)) \to
\Hom(X, S(A)[1]) \to \ldots
$$
of finite dimensional $k$-vector spaces. Namely, this sequence is
$k$-linear dual of the sequence
$$
\ldots \leftarrow
\Hom(A, X) \leftarrow
\Hom(B, X) \leftarrow
\Hom(C, X) \leftarrow
\Hom(A[1], X) \leftarrow
\ldots
$$
which is exact by Derived Categories, Lemma
\ref{derived-lemma-representable-homological}.
Next, we choose a distinguished triangle $(S(A), E, S(C), i, p, S(h))$
which is possible by axioms TR1 and TR2. We want to construct the dotted
arrow making following diagram commute
$$
\xymatrix{
S(C)[-1] \ar[r]_-{S(h[-1])} &
S(A) \ar[r]_{S(f)} &
S(B) \ar[r]_{S(g)} &
S(C) \ar[r]_{S(h)} &
S(A)[1] \\
S(C)[-1] \ar[r]^-{S(h[-1])} \ar@{=}[u] &
S(A) \ar[r]^i \ar@{=}[u] &
E \ar[r]^p \ar@{..>}[u]^\varphi &
S(C) \ar[r]^{S(h)} \ar@{=}[u] &
S(A)[1] \ar@{=}[u]
}
$$
Namely, if we have $\varphi$, then we claim for any $X$ the resulting
map $\Hom(X, E) \to \Hom(X, S(B))$ will be an isomorphism of $k$-vector
spaces. Namely, we will obtain a commutative diagram
$$
\xymatrix{
\Hom(X, S(C)[-1]) \ar[r] &
\Hom(X, S(A)) \ar[r] &
\Hom(X, S(B)) \ar[r] &
\Hom(X, S(C)) \ar[r] &
\Hom(X, S(A)[1]) \\
\Hom(X, S(C)[-1]) \ar[r] \ar@{=}[u] &
\Hom(X, S(A)) \ar[r] \ar@{=}[u] &
\Hom(X, E) \ar[r] \ar[u]^\varphi &
\Hom(X, S(C)) \ar[r] \ar@{=}[u] &
\Hom(X, S(A)[1]) \ar@{=}[u]
}
$$
with exact rows (see above) and we can apply the 5 lemma
(Homology, Lemma \ref{homology-lemma-five-lemma}) to see
that the middle arrow is an isomorphism. By the Yoneda lemma
we conclude that $\varphi$ is an isomorphism.

\medskip\noindent
To find $\varphi$, let us form the matress of Remark \ref{remark-matress}
from the distinguished triangles
$A \to B \to C \to A[1]$ and $S(A) \to E \to S(C) \to S(A)[1]$.
By Remark \ref{remark-dual-matress} the $k$-linear dual 
of this is also a matress. Using that $S$ is a Serre functor, we
see as above that this dual matress is the double complex made from
the triangles (the second of which we do not yet know to be distinguished)
$S(A) \to E \to S(C) \to S(A)[1]$ and $S(A) \to S(B) \to S(C) \to S(A)[1]$.
But the matress condition exactly tells us that the maps
$S(A) \to S(A)$ and $S(C) \to S(C)$ fit into a morphism of triangles,
i.e., that we have a $\varphi$ making the diagram above commute.
\end{proof}

\end{document}
